% Standalone problem document, generated from the adjacent README.md.
% Compile with XeLaTeX. Mathematical content is maintained in Markdown.
\documentclass[11pt,a4paper]{article}
\usepackage[fontsize=10.5pt]{scrextend}
\usepackage[margin=22mm,headheight=15pt,headsep=9mm,footskip=11mm]{geometry}
\usepackage{fontspec}
\setmainfont{texgyrepagella}[Extension=.otf,UprightFont=*-regular,BoldFont=*-bold,ItalicFont=*-italic,BoldItalicFont=*-bolditalic]
\setsansfont{texgyreheros}[Extension=.otf,UprightFont=*-regular,BoldFont=*-bold,ItalicFont=*-italic,BoldItalicFont=*-bolditalic]
\setmonofont{texgyrecursor}[Extension=.otf,UprightFont=*-regular,BoldFont=*-bold,ItalicFont=*-italic,BoldItalicFont=*-bolditalic,Scale=MatchLowercase]
\usepackage{amsmath,amssymb}
\usepackage{unicode-math}
\setmathfont{texgyrepagella-math.otf}
\usepackage{xcolor,microtype,enumitem,needspace,fancyhdr}
\usepackage{bookmark}
\definecolor{ink}{HTML}{17344B}
\definecolor{accent}{HTML}{197D80}
\definecolor{muted}{HTML}{596773}
\hypersetup{colorlinks=true,linkcolor=accent,urlcolor=accent,pdftitle={NR-01 - Exact
nonnegative rank of regular polygon slack
matrices},pdfauthor={Open Problems in Numerical Linear Algebra}}
\urlstyle{same}
\setlength{\parindent}{0pt}
\setlength{\parskip}{5pt plus 1pt minus 1pt}
\linespread{1.04}
\setlength{\emergencystretch}{3em}
\widowpenalty=10000
\clubpenalty=10000
\displaywidowpenalty=10000
\allowdisplaybreaks[1]
\setlist{nosep,leftmargin=1.5em}
\providecommand{\tightlist}{\setlength{\itemsep}{0pt}\setlength{\parskip}{0pt}}
\providecommand{\pandocbounded}[1]{#1}
\pagestyle{fancy}
\fancyhf{}
\fancyhead[L]{\sffamily\footnotesize\color{muted}OPEN PROBLEMS IN NUMERICAL LINEAR ALGEBRA}
\fancyhead[R]{\sffamily\bfseries\footnotesize\color{accent}NR-01}
\fancyfoot[L]{\parbox[b]{0.9\textwidth}{\sffamily\fontsize{8}{10}\selectfont\color{muted}Editorial ratings; a bounded search cannot certify that a problem remains unsolved.\\Literature check: 2026-09-13}}
\fancyfoot[R]{\sffamily\footnotesize\color{muted}\thepage}
\renewcommand{\headrulewidth}{0.3pt}
\renewcommand{\headrule}{\hbox to\headwidth{\color{accent}\leaders\hrule height \headrulewidth\hfill}}
\makeatletter
\renewcommand{\section}{\Needspace{5\baselineskip}\@startsection{section}{1}{0pt}{12pt plus 2pt minus 2pt}{4pt}{\sffamily\bfseries\large\color{ink}}}
\renewcommand{\subsection}{\Needspace{4\baselineskip}\@startsection{subsection}{2}{0pt}{9pt plus 1pt minus 1pt}{3pt}{\sffamily\bfseries\normalsize\color{ink}}}
\makeatother
\setcounter{secnumdepth}{0}
\begin{document}
{\sffamily\small\color{muted}Nonnegative and positive
factorizations\par}
\vspace{3pt}
{\raggedright\hyphenpenalty=10000\exhyphenpenalty=10000\sffamily\bfseries\fontsize{21}{24}\selectfont\color{ink}Exact
nonnegative rank of regular polygon slack matrices\par}
\vspace{7pt}
{\sffamily\small\textbf{Difficulty:} challenging\quad\textbf{Importance:} interesting
to the community\par}
{\sffamily\small\bfseries\color{accent}Status: Partially resolved\par}
\vspace{4pt}
{\color{accent}\hrule height 0.8pt}
\vspace{6pt}
\textbf{Rating rationale:} Challenging because matching lower bounds
must hold for every polygon size despite an explicit upper construction;
community importance connects structured NMF with optimal linear
extended formulations.

\textbf{Area:} structured nonnegative matrix factorization

\subsection{Context and notation}\label{context-and-notation}

All factorizations are over the real numbers. For
\(X\in\mathbb R_{\ge0}^{m\times n}\), define

\[\mathop{\mathrm{rank}}\nolimits_+(X)=\min\{r\ge0:X=WH,\quad
W\in\mathbb R_{\ge0}^{m\times r},\ H\in\mathbb R_{\ge0}^{r\times n}\}.\]

\subsection{Problem statement}\label{problem-statement}

For each integer \(n\ge3\), define the \(n\times n\) nonnegative matrix

\[S_n(i,j)=\cos(\pi/n)-\cos\bigl((2i+1-2j)\pi/n\bigr),
\qquad 0\le i,j< n.\]

It is the slack matrix obtained from the vertices of the regular
\(n\)-gon on the unit circle and its supporting facet inequalities.

\subsubsection{Question}\label{question}

With \(k=\lceil\log_2 n\rceil\), is it true for every \(n\ge3\) that

\[\mathop{\mathrm{rank}}\nolimits_+(S_n)=
\begin{cases}
2k-1,&2^{k-1}< n\le 2^{k-1}+2^{k-2},\\
2k,&2^{k-1}+2^{k-2}< n\le2^k?
\end{cases}\]

\subsection{References}\label{references}

Arnaud Vandaele, Nicolas Gillis, François Glineur, and Daniel Tuyttens,
\href{https://arxiv.org/html/1411.7245}{\emph{Heuristics for Exact
Nonnegative Matrix Factorization}}, Journal of Global Optimization
\textbf{65} (2016), 369--400, §6.2, Conjecture 2. Vandaele, Gillis, and
Glineur, \href{https://arxiv.org/abs/1505.08031}{\emph{On the Linear
Extension Complexity of Regular n-gons}}, Linear Algebra and its
Applications \textbf{521} (2017), 217--239, proves the corresponding
upper bound. Nicolas Gillis,
\href{https://orbi.umons.ac.be/bitstream/20.500.12907/42337/1/NMFbook_SIAM_reprint.pdf}{\emph{Nonnegative
Matrix Factorization}}, SIAM, 2020, §3.6.3.4, p.~90.

\subsection{Status check --- 2026-09-10}\label{status-check-2026-09-10}

Rechecked
\href{https://arxiv.org/html/2605.14058v2}{Baeckelant--Vandaele--Gillis
v2, Appendix A.2, Table 7}, and searched for later regular-polygon rank
resolutions. The conjectured value is attained and proved for
substantive ranges, including n from 5 through 16, but gaps remain, for
example at n=17. The upper formula remains conjecturally sharp for the
whole family; no full resolution was located.

\subsection{Reviewed partial result ---
2026-09-13}\label{reviewed-partial-result-2026-09-13}

Sidney Holden (Center for Computational Biology, Flatiron Institute,
Simons Foundation) proves in
\href{https://github.com/ajt60gaibb/OpenProblemsInNLA/blob/main/references/holden-nr01-2026-09-13/report.pdf}{Theorem
1.1} that

\[\mathop{\mathrm{rank}}\nolimits_+(S_n)=9,\qquad n=17,18,19,20.\]

The computer-assisted lower bound covers arbitrary real factorizations,
including nonsimple lifts and nongeneric projections; exact nine-term
factorizations give the matching upper bounds. Sections 2--7 supply the
rank-balance reduction, finite shadow cover, degeneration argument and
characteristic-zero justification of the modular rank certificates. The
argument also gives the lower bound nine for every size at least 21 by
the eight-facet vertex bound in Section 2; together with the known upper
construction this settles sizes 21--24.

The
\href{https://github.com/ajt60gaibb/OpenProblemsInNLA/blob/main/references/holden-nr01-2026-09-13/independent-review.md}{independent
informal Codex AI-agent audit} checks the mathematical reductions and
reproduces the complete computational pipeline. This is partial
affirmative progress: the original equality for every size remains open,
beginning with \(9\le\mathop{\mathrm{rank}}\nolimits_+(S_{25})\le10\).
In particular, sizes 25--30 and 33--42 remain unresolved by this
submission; no general formula for the remaining sizes is proved.
\textbf{NR-01 remains Partially resolved.} No Lean verification or
external human peer review is claimed.

\href{https://github.com/ajt60gaibb/OpenProblemsInNLA/blob/main/references/holden-nr01-2026-09-13/README.md}{Submission,
reproducibility and verified affiliation} ·
\href{https://github.com/ajt60gaibb/OpenProblemsInNLA/blob/main/references/holden-nr01-2026-09-13/report.tex}{Editable
manuscript}. The original ID, canonical path, target and prior
references are retained.
\end{document}
